\subsubsection{Tarjan (puntos de articulacion)}

\paragraph{Idea intuitiva.}
Encuentra los \textbf{puntos de articulacion} (cut vertices): vertices cuya remocion \textbf{desconecta} el grafo. Mismo setup que para bridges: \texttt{tin[v]}, \texttt{low[v]}. Un nodo $v$ (no raiz) es de articulacion si existe un hijo $u$ con $low[u] \ge tin[v]$. La \textbf{raiz} es de articulacion si tiene $\ge 2$ hijos en el DFS tree. La razon: para $v$ no raiz, la condicion $\ge$ indica que el subarbol de $u$ depende de $v$; quitar $v$ los desconecta. Para la raiz, perder $\ge 2$ hijos parte el DFS en $\ge 2$ componentes.

\paragraph{Propiedades clave (invariantes).}
\begin{itemize}\setlength\itemsep{0pt}
	\item $low[u] \ge tin[v]$ significa que el subarbol de $u$ no tiene back-edge ``por encima'' de $v$.
	\item La raiz es caso especial: tiene que perder mas de un hijo para que el grafo se parta.
	\item Cada nodo se evalua una sola vez.
\end{itemize}

\paragraph{Codigo clave.}
La condicion de articulacion:
\begin{verbatim}
void dfs(int v, int p) {
    used[v] = true;
    tin[v] = low[v] = timer++;
    int children = 0;
    for (int u : adj[v]) {
        if (u == p) continue;
        if (used[u]) low[v] = min(low[v], tin[u]);
        else {
            dfs(u, v);
            low[v] = min(low[v], low[u]);
            if (low[u] >= tin[v] && p != -1)  // no raiz
                isArt[v] = true;
            ++children;
        }
    }
    if (p == -1 && children > 1) isArt[v] = true;  // raiz especial
}
\end{verbatim}

\paragraph{Complejidad.}
\begin{itemize}\setlength\itemsep{0pt}
	\item $O(V + E)$.
	\item Una sola pasada de DFS.
\end{itemize}

\paragraph{Donde tocar para modificar.}
\begin{itemize}\setlength\itemsep{0pt}
	\item \textbf{Block-cut tree}: estructura bipartita con BCCs (2-vertex-connected components) y articulation points.
	\item \textbf{Verificar biconectividad}: chequear que no hay puntos de articulacion.
	\item \textbf{Componentes biconnected}: extender Tarjan para enumerar BCCs.
\end{itemize}

\paragraph{Trampas y errores frecuentes.}
\begin{itemize}\setlength\itemsep{0pt}
	\item No olvidar el caso especial de la raiz.
	\item La condicion usa \texttt{>=}, no \texttt{>}; un hijo con \texttt{low[u] == tin[v]} tambien cuenta.
	\item Multigrafos pueden tener ``falsos'' articulation points; verificar caso a caso.
\end{itemize}

\paragraph{Codigo.}
Ver \texttt{cpp.tex}, seccion \textit{Tarjan (puntos de articulacion)}.

\vspace{0.5em}
